\medterm{Jaccoud Arthropathy} A form of joint disease that can permanently or temporarily pull the fingers toward the little-finger side without the bone damage typical of rheumatoid arthritis. It is linked most often with past rheumatic fever or lupus.

\medterm{Jaccoud's Arthropathy} A deforming but usually non-damaging form of arthritis that can cause the fingers to drift toward the little-finger side. The deformity may improve when the hand is moved and is often linked to lupus or past rheumatic fever.

\medterm{Jack-In-The-Pulpit Poisoning} Jack-in-the-pulpit poisoning follows ingestion of calcium oxalate crystals from the plant, causing intense burning and swelling of the mouth, throat, and gastrointestinal tract.

\medterm{Jackknife Position} A face-down position used during some operations in which the table is bent at the hips. The buttocks are raised while the head and legs are lower, improving access to the rectal, anal, or posterior pelvic regions.

\synonyms
Kraske Position

\medterm{Jacksonian Epilepsy} A focal seizure that begins with twitching or stiffening in one body part and spreads in an orderly way to nearby body parts. Awareness may remain normal at first.

\synonyms
Jacksonian March; Focal Motor Seizure

\medterm{Jacksonian March} A clinical progression of focal motor seizure activity in which epileptic discharges propagate along contiguous areas of the primary motor cortex (precentral gyrus homunculus), resulting in clonic contractions that begin focally in one anatomical region (such as the thumb or angle of the mouth) and march sequentially to involve adjacent muscle groups of the limb or face.
\synonyms
Jacksonian Seizure; Focal Motor March

\medterm{Jacksonian Seizure} A focal motor seizure that begins in one body part and may spread in an orderly pattern to adjacent body regions; awareness may remain intact at first.

\medterm{Jacobsen Syndrome} A contiguous gene deletion disorder caused by a terminal deletion of the long arm of chromosome 11 (11q23-q24), characterized clinically by Paris-Trousseau thrombocytopenia and platelet dysfunction, trigonocephaly, facial dysmorphism, structural cardiac defects, and psychomotor delay.
\synonyms
11q Terminal Deletion Syndrome; Jacobsen Deletion Disorder

\medterm{Jadassohn's Naevus} A hairless skin patch or raised lesion present from birth, usually caused by an overgrowth of oil glands. It may change during life and rarely develops another tumor, so changing lesions need review.

\medterm{Jadassohn-Lewandowsky Syndrome} The classic form of pachyonychia congenita (type 1) caused by pathogenic mutations in the KRT6A or KRT16 keratin genes, presenting with hypertrophic subungual hyperkeratosis and thickened dystrophic nails, severe painful palmoplantar keratoderma, oral mucosal leukokeratosis, and follicular hyperkeratosis.
\synonyms
Pachyonychia Congenita Type 1; PC-1

\medterm{Jaffe-Lichtenstein Syndrome} A benign bone condition characterized by polyostotic fibrous dysplasia involving multiple skeletal sites with replacement of normal lamellar bone by immature fibro-osseous tissue, occurring in the absence of endocrine hyperfunction or the café-au-lait hyperpigmentation seen in McCune-Albright syndrome.
\synonyms
Polyostotic Fibrous Dysplasia; Jaffe-Lichtenstein Disease

\medterm{JAK-STAT Pathway} A signaling system cells use to respond to hormones and immune chemicals. It helps control growth, blood-cell production, inflammation, and immune activity; abnormal signaling can contribute to inflammatory disease or cancer.

\medterm{JAK2} A protein inside cells that passes signals from certain hormones and immune chemicals. An activating change in the JAK2 gene can make the bone marrow produce too many blood cells, as in several myeloproliferative disorders.

\medterm{Jammed Finger} A finger injury caused by sudden force, often involving a sprain, joint dislocation, fracture, or tendon injury; persistent deformity, severe pain, or inability to move requires assessment.

\medterm{Janeway Lesions} Small, usually painless red or purple spots on the palms or soles that can occur with infective endocarditis, an infection of the heart lining or valves. They result from infected particles blocking tiny blood vessels.

\medterm{Janus Kinase} One of a family of proteins inside cells that pass signals from hormones and immune chemicals. Janus kinases help control inflammation, blood-cell production, and immune responses.

\medterm{Janus Kinase Inhibitors} Medicines that block Janus kinase proteins and reduce certain immune signals. They are used for selected inflammatory and blood disorders, but can increase the risk of infections, blood clots, and other serious effects.

\medterm{Japanese Encephalitis} A viral infection spread by mosquitoes that can inflame the brain. Most infections cause no symptoms, but severe illness can cause fever, confusion, seizures, movement problems, or lasting disability.

\medterm{Jargon Aphasia} A severe fluent receptive aphasia variant (frequently associated with extensive posterior superior temporal lobe lesions in Wernicke aphasia) characterized by fluent, effortless speech filled with neologisms, phonemic paraphasias, and unintelligible syntax ('word salad') with marked impairment in comprehension.
\synonyms
Neologistic Aphasia; Word Salad Speech

\medterm{Jarisch-Herxheimer Reaction} A short-lived reaction that may occur soon after antibiotics are started for certain infections. It can cause fever, chills, headache, muscle aches, and temporarily worse symptoms.

\medterm{Jaundice} Yellowish pigmentation of the skin, sclerae, and mucous membranes caused by elevated levels of circulating bilirubin in the blood. Etiologies include pre-hepatic hemolytic processes, intra-hepatic hepatocellular injury (hepatitis, cirrhosis), or post-hepatic biliary obstruction (gallstones, strictures, neoplasms).

\synonyms
Icterus; Hyperbilirubinemic Scleral Icterus; Biliary Pigmentation

\medterm{Jaundice And Breastfeeding} Jaundice associated with breastfeeding may reflect inadequate milk intake in the early days or persistent breast-milk jaundice after other serious causes have been excluded.

\medterm{Jaundice Causes} Jaundice can result from rapid breakdown of red blood cells, inherited bilirubin problems, liver disease, medicines or toxins, or blockage of bile flow. Blood tests and imaging help identify the cause and guide treatment.

\medterm{Jaundice In Adults} Adult jaundice is yellowing of skin or eyes from elevated bilirubin due to red-cell breakdown, liver injury, impaired bile flow, or inherited metabolism. Dark urine, pale stool, abdominal pain, fever, confusion, or rapid onset requires prompt assessment.

\medterm{Jaundice In Neonates} Yellowing of a newborn's skin or eyes from elevated bilirubin. It is often physiologic, but jaundice in the first 24 hours, rapidly increasing jaundice, poor feeding, lethargy, fever, pale stools, or dark urine requires urgent evaluation.

\medterm{Jaundice In Newborns} Yellowing of a newborn’s skin or eyes caused by increased bilirubin. It is often temporary, but very high or rapidly rising levels can injure the brain and require urgent treatment.

\medterm{Jaw} The upper and lower bony structures of the mouth—the maxilla and mandible—which support the teeth and contribute to chewing, speech, facial shape, and the oral airway.

\medterm{Jaw Fracture} A break in the lower jawbone or upper jawbone, usually after a blow to the face. Signs may include pain, swelling, loose teeth, numbness, a changed bite, or difficulty opening the mouth. Severe breaks can block the airway, especially when the lower jaw is broken in the front on both sides. Treatment may require realignment, wiring, plates, or surgery.

\medterm{Jaw Jerk Reflex} A brief closing movement of the jaw when the chin is tapped with the mouth slightly open. An unusually strong response may suggest a problem in the brain's motor pathways.

\medterm{Jaw Neoplasm} An abnormal growth in the upper or lower jawbone or nearby tissues. It may be a cyst, a benign tumor, or cancer and can cause swelling, pain, loose teeth, numbness, or bone destruction.

\medterm{Jaw Pain} Pain in the jaw, jaw joint, or chewing muscles caused by temporomandibular disorders, dental disease, infection, injury,
or less commonly heart or nerve disease. Pain with chest discomfort, sweating, breathlessness, or exertion
requires urgent assessment; persistent pain or restricted jaw movement requires dental or medical review.

\medterm{Jaw Tumors And Cysts} Abnormal growths or fluid-filled sacs developing in the jawbone or nearby tissues. They may be harmless or cancerous and can cause swelling, pain, tooth movement, numbness, or bone destruction.

\synonyms
Odontogenic Tumors And Cysts

\medterm{Jaw-Winking Syndrome} A congenital synkinetic neuro-ophthalmic condition (Marcus Gunn phenomenon) characterized by unilateral congenital ptosis that transiently resolves or retracts upward with movement of the lower jaw during mastication, sucking, or lateral pterygoid contraction, caused by anomalous congenital innervation of the levator palpebrae superioris by the motor branch of the trigeminal nerve (CN V3).
\synonyms
Marcus Gunn Jaw-Winking Phenomenon; Trigemino-Oculomotor Synkinesis

\medterm{Jehovah's Witness} A member of a Christian religious group whose beliefs may lead them to decline transfusions of whole blood and certain blood components. Individual choices about blood products, procedures, and treatment vary, so clinicians should discuss the patient's informed preferences and document them carefully.

\medterm{Jejunal Atresia} A condition present at birth in which part of the jejunum, a section of the small intestine, is closed or missing. It causes bowel blockage in a newborn, with green vomiting, a swollen abdomen, and failure to pass the first stool.

\medterm{Jejunal Fistula} An abnormal connection involving the jejunum and another organ, skin, or body surface. It may cause intestinal contents to leak, infection, malnutrition, dehydration, or sepsis and requires assessment of its cause and output.

\medterm{Jejunal Hemorrhage} Bleeding from the jejunum, the middle segment of the small intestine, which may cause maroon or black stool, anemia, abdominal pain, or shock depending on its rate and cause.

\medterm{Jejunal Neoplasm} An abnormal growth in the jejunum, the middle part of the small intestine. It may be benign or cancerous and can cause pain, bleeding, anemia, vomiting, blockage, or weight loss.

\medterm{Jejunal Stenosis} Narrowing of the jejunal lumen that can impede intestinal contents and cause cramping, vomiting, distension, malabsorption, or obstruction; causes include Crohn disease, adhesions, tumors, and ischemic injury.

\medterm{Jejunal Ulcer} An ulcerated break in the lining of the jejunum, caused by inflammation, infection, ischemia, medicines, tumor, or anastomotic disease and potentially producing pain, bleeding, or perforation.

\medterm{Jejunitis} Jejunitis is inflammation of the jejunum, which may result from infection, inflammatory bowel disease, ischemia, medication, or other intestinal disorders.

\medterm{Jejunoileal Atresia} A birth defect in which part of the small intestine is completely closed or missing, so milk and digestive fluid cannot pass. Newborns develop green vomiting, a swollen abdomen, and failure to pass stool. Imaging confirms the blockage, and surgery is needed to restore intestinal passage.

\synonyms
Jejunal Atresia; Ileal Atresia; Small Bowel Atresia

\medterm{Jejunostomy} A surgically created opening into the jejunum, often used for placement of a feeding tube when nutrition cannot safely be given through the stomach.

\medterm{Jejunostomy Feeding Tube} A jejunostomy feeding tube delivers nutrition directly into the jejunum when oral or gastric feeding is unsafe or inadequate, with attention to formula, rate, position, and complications.

\medterm{Jejunum} The middle part of the small intestine. It absorbs much of the nutrients from food before the remaining contents pass into the ileum.

\medterm{Jellyfish Sting} A painful skin reaction caused when a jellyfish tentacle releases venom into the skin. It may cause burning, redness, and raised lines; severe exposure can cause vomiting, breathing trouble, fainting, or widespread symptoms.

\synonyms
Jellyfish Envenomation; Sea Sting

\medterm{Jellyfish Stings} Alternate name; see \textit{Jellyfish stings}.
\medterm{Jendrassik Maneuver} A clinical examination technique used to reinforce and elicit sluggish deep tendon reflexes (such as the patellar tendon reflex) by having the patient hook their flexed fingers together and pull forcefully apart, which increases spinal motoneuron pool excitability by reducing descending central inhibition.
\synonyms
Jendrassik Reinforcement; Reflex Reinforcement Maneuver

\medterm{Jendrassik's Maneuver} A technique used during a reflex examination in which a person pulls the hands apart or tenses another muscle group. The distraction can make a knee or other tendon reflex easier to observe.

\medterm{Jerk Nystagmus} An involuntary, rhythmic oscillation of the eyes characterized by a slow drift in one direction followed by a rapid corrective saccadic movement in the opposite direction. The direction of the nystagmus is designated by the fast phase, and it may arise from vestibular, cerebellar, or brainstem pathology.

\medterm{Jerusalem Cherry Poisoning} Jerusalem cherry poisoning follows ingestion of Solanum pseudocapsicum berries, which contain toxic compounds that may cause gastrointestinal symptoms, confusion, weakness, or cardiac effects.

\medterm{Jervell and Lange-Nielsen Syndrome} A rare inherited condition causing severe hearing loss from birth and an abnormal heart rhythm. The rhythm problem can cause fainting, seizures, or sudden death, so urgent specialist care is important.

\medterm{Jet Lag} Temporary circadian rhythm disruption after rapid travel across time zones, causing sleep disturbance, fatigue, impaired concentration, and sometimes gastrointestinal symptoms.

\synonyms
Desynchronosis; Time Zone Change Syndrome

\medterm{Jet Lag Prevention} Jet lag prevention uses timed light exposure, sleep adjustment, hydration, and carefully planned activity or melatonin strategies to shift the circadian rhythm across time zones.

\medterm{Jeune Syndrome} An autosomal recessive skeletal ciliopathy caused by mutations in intraflagellar transport genes, characterized by a severely constricted, narrow bell-shaped thorax, short-limbed micromelic skeletal dysplasia, polydactyly, pulmonary hypoplasia, and progressive cystic renal and hepatic disease.
\synonyms
Asphyxiating Thoracic Dystrophy; ATD

\medterm{Jewelry Cleaner Poisoning} Jewelry cleaner poisoning can injure the mouth, esophagus, stomach, lungs, or eyes depending on the product and exposure route; immediate poison-control guidance is appropriate.

\medterm{Jimsonweed Poisoning} Jimsonweed poisoning causes anticholinergic toxicity, including dilated pupils, dry mouth, rapid heart rate, urinary retention, agitation, delirium, seizures, or coma.

\medterm{Job Syndrome} A rare inherited immune disorder causing eczema, repeated skin and lung infections, and very high levels of an allergy-related antibody. Dental, bone, and facial differences may also occur.

\medterm{Jock Itch} A fungal infection of the groin and inner thighs, also called tinea cruris. It causes an itchy, scaly, ring-shaped or spreading
rash, often worsened by warmth and moisture. Pain, fever, extensive involvement, or failure to improve with
appropriate care requires clinical evaluation.

\synonyms
Tinea Cruris; Ringworm of the Groin

\medterm{Jod-Basedow Phenomenon} Overactive thyroid function triggered by an excess of iodine, usually in a person with thyroid nodules or other thyroid tissue that works independently of normal control. It may cause weight loss, tremor, sweating, or a fast heartbeat.

\medterm{Joffroy Sign} A physical examination sign of Graves disease and hyperthyroidism characterized by the absence of normal forehead wrinkling and furrowing when the patient looks upward while keeping the head in a fixed forward position.
\synonyms
Joffroy's Sign; Thyrotoxic Forehead Sign

\medterm{Johanson-Blizzard Syndrome} A rare autosomal recessive congenital malformation syndrome caused by mutations in the UBR1 gene, characterized by exocrine pancreatic insufficiency, hypoplasia or aplasia of the nasal alae ('beaked nose'), scalp vertex defects, oligodontia, sensorineural hearing loss, and growth retardation.
\synonyms
JBS; Pancreatic Insufficiency-Alar Hypoplasia Syndrome

\medterm{Joint} The place where two or more bones meet. Joints may be fixed, slightly movable, or freely movable, such as the knee, hip, or shoulder.

\medterm{Joint Aspiration} A needle procedure that removes fluid from a joint for testing or symptom relief, helping assess infection, bleeding, crystals, or inflammation.
\medterm{Joint Capsule} A tough envelope surrounding a freely movable joint. Its inner lining makes slippery fluid that helps the joint move, while the capsule helps hold the bones together and provides stability.

\medterm{Joint Contracture} Permanent or long-lasting limitation of joint movement caused by tightening or scarring of muscles, tendons, ligaments, the joint capsule, skin, or nearby tissues. It may follow immobility, injury, burns, inflammation, or nerve disease.

\medterm{Joint Diseases} Conditions that affect joints, including osteoarthritis, inflammatory arthritis, gout, infection, injury, and tumors. They may cause pain, swelling, stiffness, deformity, or reduced movement; treatment depends on the cause.

\medterm{Joint Effusion} Extra fluid inside a joint, often causing swelling, stiffness, or restricted movement. Causes include injury, infection, inflammatory arthritis, bleeding, and crystal disease; a clinician may remove fluid for testing.

\medterm{Joint Fluid Culture} Joint fluid culture tests synovial fluid for bacterial, fungal, or other organisms when septic arthritis is suspected; results are interpreted with cell count, crystals, and clinical findings.

\medterm{Joint Fluid Gram Stain} A quick laboratory test that stains fluid removed from a swollen joint to look for bacteria. It can support the evaluation of a serious joint infection but a negative result does not rule infection out.

\medterm{Joint Inflammation} Alternate name; see \textit{Arthritis}.

\medterm{Joint Mouse} A loose piece of cartilage, bone, or other tissue inside a joint that can cause catching, locking, pain, or swelling.

\medterm{Joint Pain} Pain or aching in one or more joints or in the tissues around them; the medical term is arthralgia. It is a symptom rather than a single disease and may result from arthritis, inflammation, injury, overuse, infection, crystals such as those in gout, or problems in nearby muscles, tendons, or bursae. It may occur with stiffness, swelling, warmth, redness, tenderness, or reduced movement. The number of joints affected, the timing, and associated symptoms help identify the cause.

\medterm{Joint Prosthesis} An artificial joint or implant designed to replace a damaged or diseased biological
joint, restoring mobility and relieving chronic pain.

\medterm{Joint Replacement} Surgery that removes damaged joint surfaces and replaces them with artificial parts to reduce pain and improve movement. It is commonly performed for severe arthritis or joint destruction and requires rehabilitation.

\medterm{Joint Replacement Rehabilitation} Exercise, education, and daily-activity training after joint replacement, helping restore safe movement, strength, balance, confidence, and independence.

\medterm{Joint Stiffness} Reduced ease or range of joint movement, often worse after rest. Causes include arthritis, injury,
inflammation, prolonged immobility, or muscle and tendon disorders. Sudden swelling, redness,
fever, severe pain, or inability to use the joint requires assessment.

\medterm{Joint Swelling} Enlargement of a joint caused by excess synovial fluid, blood, inflammation, infection,
trauma, or structural disease. It may be accompanied by pain, warmth, stiffness, or
restricted movement and can indicate inflammatory or septic arthritis.

\medterm{Joint X-Ray} An X-ray picture of a joint used to look for broken bones, abnormal alignment, arthritis-related wear, bone loss, deposits, or other changes. It uses a small amount of radiation and does not show all soft-tissue injuries.

\medterm{Jones Criteria} A standardized set of diagnostic guidelines formulated by the American Heart Association to establish the diagnosis of acute rheumatic fever following group A streptococcal pharyngitis, requiring evidence of antecedent streptococcal infection along with two major criteria (carditis, migratory polyarthritis, Sydenham chorea, erythema marginatum, subcutaneous nodules) or one major and two minor criteria.
\synonyms
Revised Jones Criteria; AHA Rheumatic Fever Criteria

\medterm{Jones Fracture} An acute, traumatic, extra-articular transverse fracture of the fifth metatarsal bone located at the metaphyseal-diaphyseal junction (approximately 1.5 cm distal to the tuberosity), prone to delayed union and nonunion due to poor vascular supply in this watershed zone.
\synonyms
Fifth Metatarsal Watershed Fracture; Metaphyseal-Diaphyseal Fifth Metatarsal Fracture

\medterm{Joubert Syndrome} A rare inherited disorder of brain development that affects breathing control, balance, eye movements, and muscle tone. Babies may have pauses or rapid changes in breathing, poor muscle tone, delayed development, and abnormal eye movements. Some also have kidney, liver, or retinal disease, or extra fingers or toes. Brain MRI helps confirm the diagnosis.

\synonyms
JSRD; Joubert-Related Disorder

\medterm{Jugular} Relating to the neck, especially the jugular veins that carry blood from the head and neck back to the heart.

\medterm{Jugular Foramen} An opening at the base of the skull through which a major neck vein and several nerves pass. Disease in this area can cause difficulty swallowing, hoarseness, shoulder weakness, hearing symptoms, or other neurologic problems.

\medterm{Jugular Foramen Syndrome} A clinical neurological syndrome resulting from compressive or infiltrative lesions (such as glomus jugulare tumors, schwannomas, or skull base metastases) affecting the structures traversing the jugular foramen, characterized by ipsilateral paralysis of the glossopharyngeal (CN IX), vagus (CN X), and accessory (CN XI) cranial nerves, leading to dysphagia, vocal cord paralysis, and trapezius weakness.
\synonyms
Vernet Syndrome; Posterior Lacerated Foramen Syndrome

\medterm{Jugular Vein} A major vein that carries blood from the head, face, and neck back to the heart. Each side has an internal jugular vein and an external jugular vein.

\medterm{Jugular Venous Distension} Pathological engorgement and visible protrusion of the internal or external jugular veins in the neck above the level of the clavicle with the patient positioned at 45 degrees, providing a direct physical reflection of elevated right atrial and central venous pressure, classic for right heart failure, cardiac tamponade, and constrictive pericarditis.
\synonyms
JVD; Raised Jugular Venous Pressure

\medterm{Jugular Venous Pressure} The visible height of blood in the neck veins. It helps clinicians estimate pressure in the right side of the heart and assess fluid buildup.

\synonyms
JVP

\medterm{Jugulodigastric Node} A prominent deep cervical lymph node situated at the junction where the posterior belly of the digastric muscle crosses the internal jugular vein, receiving primary lymphatic drainage from the palatine tonsils, base of tongue, and posterior pharynx, frequently enlarged in acute pharyngotonsillitis and head and neck neoplasms.
\synonyms
Tonsillar Lymph Node; Subdigastric Node

\medterm{Jumper's Knee} Alternate name; see \textit{Patellar Tendinitis}.

\medterm{Junctional Escape Rhythm} A physiological backup cardiac rhythm originating from pacemaking cells within the atrioventricular (AV) junction and bundle of His at an intrinsic rate of 40 to 60 beats per minute, characterized by narrow QRS complexes with absent, inverted, or retrograde P waves, arising when sinoatrial node pacemaking fails or during complete AV nodal block.
\synonyms
AV Junctional Rhythm; Nodal Escape Rhythm

\medterm{Junctional Nevus} A usually benign mole formed by melanocytes at the junction between the epidermis and dermis; it is often flat and should be assessed if it changes in size, shape, color, or symptoms.

\medterm{Junctional Rhythm} A heart rhythm that starts near the connection between the upper and lower chambers instead of the heart’s usual pacemaker. It may be slow or regular and can occur when the usual pacemaker is suppressed or blocked.

\medterm{Junctional Tachycardia} A supraventricular arrhythmia arising from enhanced automaticity or triggered activity within the AV node or bundle of His, characterized by a rapid heart rate (greater than 100 beats per minute), regular narrow QRS complexes, and retrograde or dissociated atrial activation, frequently associated with digitalis toxicity, myocardial ischemia, or pediatric cardiac surgery.
\synonyms
Nonparoxysmal Junctional Tachycardia; Accelerated Junctional Rhythm

\medterm{Jungian Analysis} Jungian analysis is a psychotherapy approach based on Carl Jung’s theories, using exploration of unconscious processes, symbols, dreams, and personality development.

\medterm{Juvenile Absence Epilepsy} An idiopathic generalized epilepsy syndrome with onset in adolescence (typically between 9 and 13 years of age), characterized by absence seizures occurring less frequently than in childhood absence epilepsy, accompanied on EEG by generalized 3 to 4 Hz spike-and-wave discharges, frequently accompanied by generalized tonic-clonic seizures.
\synonyms
JAE; Adolescent Absence Epilepsy

\medterm{Juvenile Angiofibroma} A noncancerous but highly blood-vessel-rich tumor that usually develops near the back of the nasal cavity in adolescent boys. It may cause worsening nasal blockage, repeated nosebleeds, or facial swelling and should be assessed by specialists.

\medterm{Juvenile Arthritis} A general term for inflammatory joint disease beginning in childhood. It may cause persistent swelling, pain, stiffness, reduced movement, fever, rash, or eye inflammation; the subtype determines monitoring and treatment.

\synonyms
Juvenile Idiopathic Arthritis; JIA; Juvenile Rheumatoid Arthritis

\medterm{Juvenile Astrocytoma} A tumor arising from support cells in a child’s or adolescent’s brain or spinal cord. Many grow slowly, but symptoms depend on location and may include headache, seizures, vision changes, weakness, balance problems, or hormone disturbances.

\medterm{Juvenile Dermatomyositis} An autoimmune illness in children that weakens muscles, especially around the hips and shoulders, and causes a purple rash around the eyes or raised patches over the knuckles. It may also affect swallowing, breathing, or other organs.

\medterm{Juvenile Diabetes} Alternate historical term; see \textit{Type 1 Diabetes}.

\medterm{Juvenile Hypermobile Spectrum Disorder} A hypermobility spectrum disorder beginning in childhood or adolescence, with excessive joint movement plus pain, instability, soft-tissue injury, fatigue, or related symptoms.
\medterm{Juvenile Hypermobility Syndrome} A pattern of unusually flexible joints in a child or adolescent that causes pain, joint instability, tiredness, or repeated sprains. Assessment checks for other connective-tissue disorders; treatment may include strengthening, activity guidance, and joint protection.

\medterm{Juvenile Idiopathic Arthritis} A group of inflammatory joint diseases that begin before age 16 and last at least six weeks without another identified cause. Children may have swelling, stiffness, pain, fever, rash, or eye inflammation, depending on the subtype.

\medterm{Juvenile Myoclonic Epilepsy} A generalized epilepsy syndrome that typically begins in adolescence and causes myoclonic jerks, often after awakening, with possible generalized tonic-clonic or absence seizures.

\medterm{Juvenile Nasopharyngeal Angiofibroma} A benign, highly vascular, unencapsulated mesenchymal tumor occurring almost exclusively in adolescent males, originating in the posterolateral wall of the nasal cavity near the sphenopalatine foramen, presenting clinically with recurrent painless epistaxis, nasal obstruction, and progressive facial swelling.
\synonyms
JNA; Nasopharyngeal Fibroangioma

\medterm{Juvenile Osteoporosis} Low bone strength during childhood or adolescence that increases the risk of fractures. It may occur on its own or result from chronic illness, poor nutrition, hormone disorders, limited mobility, or medicines.

\medterm{Juvenile Polyposis Syndrome} An inherited condition causing many noncancerous growths in the stomach or intestines. The growths can bleed, cause anemia, and increase the risk of bowel and stomach cancer, so regular screening is needed.

\medterm{Juvenile Rheumatoid Arthritis} Alternate historical term; see \textit{Juvenile Idiopathic Arthritis}.

\medterm{Juvenile Schizophrenia} Alternate name; see \textit{Childhood Schizophrenia}.

\medterm{Juvenile Xanthogranuloma} A rare, usually harmless growth of immune cells in the skin of an infant or young child, appearing as a firm red, yellow, or brown bump. It often fades without treatment, but rarely affects the eye or internal organs.

\medterm{Juxta-Articular} Located next to a joint. The term may describe bone, soft tissue, swelling, a cyst, or a deposit beside a joint rather than within the joint itself.

\medterm{Juxtaglomerular Apparatus} A group of specialized kidney cells that senses blood flow and salt levels and releases renin. It helps regulate blood pressure, blood volume, and the body’s salt balance.

\medterm{Juxtaglomerular Cell Tumor} A rare, benign renal neoplasm arising from the specialized renin-secreting myoepithelial juxtaglomerular cells of the afferent arteriole, characterized by marked hyperreninemia, secondary hyperaldosteronism, severe hypertension, and hypokalemia.
\synonyms
Reninoma; Juxtaglomerular Renin-Secreting Tumor

\medterm{Juxtavesical} Next to the urinary bladder. The term often describes the lower part of a ureter as it approaches the bladder, where a passing stone may become lodged and cause side pain, urinary frequency, or painful urination.
